%%% LaTeX Template: Designer's CV
%%%
%%% Source: http://www.howtotex.com/
%%% Feel free to distribute this template, but please keep the referal to HowToTeX.com.
%%% Date: March 2012


%%%%%%%%%%%%%%%%%%%%%%%%%%%%%%%%%%%%%
% Document properties and packages
%%%%%%%%%%%%%%%%%%%%%%%%%%%%%%%%%%%%%
\documentclass[letterpaper,12pt,final]{memoir}

% misc
\renewcommand{\familydefault}{bch}	% font
\pagestyle{empty}					% no pagenumbering
\setlength{\parindent}{0pt}			% no paragraph indentation


% required packages (add your own)
\usepackage{flowfram}										% column layout
\usepackage[top=1cm,left=0.2cm,right=1cm,bottom=0.8cm]{geometry}% margins
\usepackage{graphicx}
\usepackage{epstopdf} 									% figures
\usepackage{url}											% URLs
\usepackage[usenames,dvipsnames]{xcolor}					% color
\usepackage{multicol}										% columns env.
	\setlength{\multicolsep}{1pt}
\usepackage{paralist}										% compact lists
\usepackage{tikz}
\usepackage{textcomp}
\usepackage{amsmath,amsfonts}
%\usepackage[none]{hyphenat}
\usepackage{geometry}
\geometry{
	%letterpaper,
	%total={215.9mm,279.4mm},
	left=8mm,
	right=18mm,
	%top=7mm,
	%bottom=7mm,
}

%\usepackage{anyfontsize}
\usepackage[hidelinks]{hyperref}
\usepackage[pscoord]{eso-pic}% The zero point of the coordinate systemis the lower left corner of the page (the default).
\usepackage{xcolor}

%%%%%%%%%%%%%%%%%%%%%%%%%%%%%%%%%%%%%
% Create column layout
%%%%%%%%%%%%%%%%%%%%%%%%%%%%%%%%%%%%%
% define length commands
\righthyphenmin =62
\lefthyphenmin =62

\setlength{\vcolumnsep}{\baselineskip}
\setlength{\columnsep}{\vcolumnsep}

% frame setup (flowfram package)
% left frame
\newflowframe{0.265\textwidth}{\textheight}{0pt}{0pt}[left]
	\newlength{\LeftMainSep}
	\setlength{\LeftMainSep}{0.21\textwidth}
	\addtolength{\LeftMainSep}{4\columnsep}

%\renewcommand{\ffvrule}[3]{%
%	\hfill
%	\tikz{\draw[loosely dotted,color=Plum,line width=1.5pt,yshift=0](0,0) -- (0,0);}
%	\hfill\mbox{}}
%\insertvrule{flow}{1}{flow}{2}

% right frame
\newflowframe{0.7\textwidth}{\textheight}{\LeftMainSep}{0pt}[main01]

% horizontal rule between frames (using TikZ)
\renewcommand{\ffvrule}[3]{%
\hfill
\tikz{\draw[loosely dotted,color=Plum,line width=1.5pt,yshift=-#1](0,0) -- (0,#3);}
\hfill\mbox{}}
\insertvrule{flow}{1}{flow}{2}

%%%%%%%%%%%%%%%%%%%%%%%%%%%%%%%%%%%%%
% define macros (for convience)
%%%%%%%%%%%%%%%%%%%%%%%%%%%%%%%%%%%%%
\newcommand{\Sep}{\vspace{1.5em}}
\newcommand{\SmallSep}{\vspace{0.5em}}

\newenvironment{Objective}
	{\ignorespaces\textbf{\color{Plum} Objective}}
	{\Sep\ignorespacesafterend}
	
\newcommand{\CVSection}[1]
	{\Large\textbf{\textsc{{#1}}}\par
	\SmallSep\normalsize\normalfont}

\newcommand{\CVItem}[1]
	{\textsc{\color{Plum} #1}}

\newcommand{\CVItemSC}[1]
{{\textsc{\color{Plum} #1}}}

\newcommand{\placetextbox}[3]{% \placetextbox{<horizontal pos>}{<vertical pos>}{<stuff>}
	\setbox0=\hbox{#3}% Put <stuff> in a box
	\AddToShipoutPictureFG*{% Add <stuff> to current page foreground
		\put(\LenToUnit{#1\paperwidth},\LenToUnit{#2\paperheight}){\vtop{{\null}\makebox[0pt][c]{#3}}}%
	}%
}%

\newcommand{\plumBox}[1]
{\space\space\fbox{\tiny \textbf{{#1}}}}

\newcommand{\RightAlignedInlineText}[1]
{{\footnotesize \color{Plum}  \hfill [#1]}}

%%%%%%%%%%%%%%%%%%%%%%%%%%%%%%%%%%%%%
% Begin document
%%%%%%%%%%%%%%%%%%%%%%%%%%%%%%%%%%%%%
\begin{document}

% Left frame
%%%%%%%%%%%%%%%%%%%%
%\begin{figure} 
%	\hfill
%	\includegraphics[width=0.8\columnwidth]{kkk}
%	\vspace{-5cm}
%\end{figure}

\begin{flushright} 
	\footnotesize
	\SmallSep
	\CVItemSC{Address}\\%{\bfseries{\color{Plum}{Address}}}\\
	
	105 North-west 103rd St.\\
	Seattle, WA - 98177\\
	\SmallSep
	\CVItemSC{E-Mail}\\%{\bfseries{\color{Plum}{E-Mail}}}\\
	\href{mailto:vicky.p.katara@gmail.com}{\texttt{vicky.p.katara@gmail.com}}\\
	\SmallSep
	\CVItemSC{Mobile}\\%{\bfseries{\color{Plum}{Cellphone}}}\\
	\href{tel:+19842158067}{+1 \(984\) 215 8067}\\
	\SmallSep
		%{\bfseries{\color{Plum}{Website}}}\\
	%\href{http://vickykatara.orgfree.com/}{vickykatara.orgfree.com}\\
	%\SmallSep
	%\CVItemSC{GitHub Repo}\\%{\bfseries{\color{Plum}{Cellphone}}}\\
	%\href{https://github.com/vicky-katara}{\texttt{github.com/vicky-katara}}\\
	\SmallSep
	\CVItemSC{Web}\\%{\bfseries{\color{Plum}{Cellphone}}}\\
	\href{https://vicky-katara.github.io/}{\texttt{vicky-katara.github.io}}\\
	\SmallSep
		\tikz{\draw[loosely dotted,color=Plum,line width=1.5pt](6.2,0) -- (1,0);}
%\end{flushright}\normalsize
%\SmallSep
\CVSection{Technical Skills}
{\footnotesize Java, Spring}\\
{\footnotesize AWS micro-services}\\
{\footnotesize Full-Stack Dev}\\
{\footnotesize JavaScript, Web-Dev}\\
{\footnotesize Dev-Ops with CI/CD}\\
\SmallSep
%\end{flushright}\normalsize
\SmallSep
\CVSection{Education}


\CVItem{Masters in \\Computer Science}\\
\begin{footnotesize}
	North Carolina State University
	2015-2016 / G.P.A.: 4.0
\end{footnotesize}
\SmallSep

\CVItem{Bachelors in \\Computer Engineering}\\
\begin{footnotesize}
    University of Mumbai
	2009-2013
\end{footnotesize}
\SmallSep

%\Sep\\
\end{flushright}\normalsize
\tikz{\draw[loosely dotted,color=Plum,line width=1.5pt](6.2,0) -- (1,0);}
\SmallSep
%\begin{figure}
	%\hfill
	%\vspace{-0.3cm}
	\hspace{0.7cm}
		\vspace{-0.35cm}
	\textsc{\bfseries{\color{Plum}{Save My Contact}}}\\\\
	\vspace{+0.45cm}
	%\hspace{-0.4cm}
	\includegraphics[width=0.8\columnwidth]{qrcode_2.eps}
%\end{figure}
\framebreak

% Right frame
%%%%%%%%%%%%%%%%%%%%
\placetextbox{0.90}{0.995}{\color{Plum}\tiny{Last Updated: Mar 24, 2024}}%

\Huge\bfseries \textsc{\color{Plum} Vicky Katara}
{\normalsize \bfseries  Senior Software Engineer }
\SmallSep

\normalsize\normalfont

% About me
%\begin{Objective}
\begin{footnotesize}
I am a seasoned senior software engineer with a strong background in Full Stack and backend micro-service development. I've lead teams of engineers to build several high value, sustainable, products and services in the AWS  and Amazon Retail product families.
\end{footnotesize}
%\end{Objective}

% Experience
\CVSection{Work Experience}

\CVItem{2021 \textendash \space Present, Senior Software Engineer}\\
\SmallSep
AWS Computer Vision (Rekognition) \& Amazon Lists\\
{\footnotesize Helped launch AWS Rekognition's (CV) custom moderation offering. Owned and lead operational excellence initiatives.\\
Tech Lead with Amazon Lists: Lead the team to launch several successful products including enhanced hearts experience add-to-list redesign, migration off of the legacy Gurupa framework (7 years in the making). Lead operational excellence for the team, including roadmap planning. Assisted with talent-reviews, mentoring, on-boarding and hiring activities}

\CVItem{2017 \textendash \space 2021, Software Engineer, Amazon}\\
\SmallSep
Amazon Lists\\
{\footnotesize Project Lead with the Amazon Lists team to deliver commercially successful products (Hearts on Amazon Products, Collaborative Lists, List-Creation Redesign. Operational leader responsible for several Ops processes within the team.}

\CVItem{May 2016 \textendash \space Aug 2016, Software Dev Intern, Amazon}\\
\SmallSep
Amazon Lists\\
{\footnotesize Built and released REST APIs for Amazon Lists. These were used for several internal migration from legacy frameworks.}
\SmallSep\\
\CVItem{Jan 2014 \textendash \space Jun 2015, Technology Analyst, Deloitte Consulting}\\
\SmallSep
Business Intelligence - Data Warehousing\\
{\footnotesize Quality engineering projects for Cellular enterprise clients}
\SmallSep

% Projects - Start
\CVSection{Projects}
\CVItem{AWS Rekognition - Custom Moderation}  \RightAlignedInlineText{2023}\\
{\footnotesize Helped launch a revolutionary enhancement for AWS Rekognition's Moderation API which allows customers to train a personalized adapter using their data, and use this to run inference.\\
	\emph{Technologies Used:} Java, Python, AWS CDK-Typescript, etc. }%Project-4 End
\SmallSep\\
%Project-new1 Start
\CVItem{Hearts and Enhanced Hearts for Mobile} \RightAlignedInlineText{2018 / 2021}\\
{\footnotesize Lead the development and launch of Amazon's Heart feature. This is used globally  on all Amazon product detail pages. The initial launch was extended by offering the customer the ability to save items to many of their lists.\\\color{Black}
\emph{Technologies Used:} Java, Spring, JavaScript, AWS Elasticsearch (etc.)}%Project-4 End
\SmallSep\\
%Project-new1 Start
\CVItem{Amazon WishLists - Gurupa Migration} \RightAlignedInlineText{2015 - 2022}\\
{\footnotesize Participated and eventually lead the Amazon WishList team's migration off of the Gurupa framework. This included the migrtion of more than 100 APIs used across many Amazon Retail Web-pages to modern microservices.\\\color{Black}
	\emph{Technologies Used:} Java, Perl-Mason, JavaScript, AWS DynamDB, etc.}%Project-4 End
\SmallSep\\
%Project-new1 Start
%Project-new1 Start
%Project-5 Start
%\CVItem{Development of Mini-Tennis Game, Mini-Project for Computer Graphics}\\
%{\footnotesize This involved the design and development of 2D Mini-Tennis game.\\ \emph{Technologies Used:} GCC}\\%Project-5 End
%\Sep 
% Projects - End
%%%%%%%%%%%%%%%%%%%%%%%%%%%%%%%%%%%%%
% End document
%%%%%%%%%%%%%%%%%%%%%%%%%%%%%%%%%%%%%
\end{document}